% Options for packages loaded elsewhere
% Options for packages loaded elsewhere
\PassOptionsToPackage{unicode}{hyperref}
\PassOptionsToPackage{hyphens}{url}
\PassOptionsToPackage{dvipsnames,svgnames,x11names}{xcolor}
%
\documentclass[
  11pt,
  a4paper,
  DIV=11,
  numbers=noendperiod]{scrartcl}
\usepackage{xcolor}
\usepackage[top=3cm, bottom=2cm, left=4cm, right=2cm]{geometry}
\usepackage{amsmath,amssymb}
\setcounter{secnumdepth}{-\maxdimen} % remove section numbering
\usepackage{iftex}
\ifPDFTeX
  \usepackage[T1]{fontenc}
  \usepackage[utf8]{inputenc}
  \usepackage{textcomp} % provide euro and other symbols
\else % if luatex or xetex
  \usepackage{unicode-math} % this also loads fontspec
  \defaultfontfeatures{Scale=MatchLowercase}
  \defaultfontfeatures[\rmfamily]{Ligatures=TeX,Scale=1}
\fi
\usepackage{lmodern}
\ifPDFTeX\else
  % xetex/luatex font selection
  \setmainfont[]{Avenir Next}
\fi
% Use upquote if available, for straight quotes in verbatim environments
\IfFileExists{upquote.sty}{\usepackage{upquote}}{}
\IfFileExists{microtype.sty}{% use microtype if available
  \usepackage[]{microtype}
  \UseMicrotypeSet[protrusion]{basicmath} % disable protrusion for tt fonts
}{}
\makeatletter
\@ifundefined{KOMAClassName}{% if non-KOMA class
  \IfFileExists{parskip.sty}{%
    \usepackage{parskip}
  }{% else
    \setlength{\parindent}{0pt}
    \setlength{\parskip}{6pt plus 2pt minus 1pt}}
}{% if KOMA class
  \KOMAoptions{parskip=half}}
\makeatother
% Make \paragraph and \subparagraph free-standing
\makeatletter
\ifx\paragraph\undefined\else
  \let\oldparagraph\paragraph
  \renewcommand{\paragraph}{
    \@ifstar
      \xxxParagraphStar
      \xxxParagraphNoStar
  }
  \newcommand{\xxxParagraphStar}[1]{\oldparagraph*{#1}\mbox{}}
  \newcommand{\xxxParagraphNoStar}[1]{\oldparagraph{#1}\mbox{}}
\fi
\ifx\subparagraph\undefined\else
  \let\oldsubparagraph\subparagraph
  \renewcommand{\subparagraph}{
    \@ifstar
      \xxxSubParagraphStar
      \xxxSubParagraphNoStar
  }
  \newcommand{\xxxSubParagraphStar}[1]{\oldsubparagraph*{#1}\mbox{}}
  \newcommand{\xxxSubParagraphNoStar}[1]{\oldsubparagraph{#1}\mbox{}}
\fi
\makeatother

\usepackage{color}
\usepackage{fancyvrb}
\newcommand{\VerbBar}{|}
\newcommand{\VERB}{\Verb[commandchars=\\\{\}]}
\DefineVerbatimEnvironment{Highlighting}{Verbatim}{commandchars=\\\{\}}
% Add ',fontsize=\small' for more characters per line
\usepackage{framed}
\definecolor{shadecolor}{RGB}{241,243,245}
\newenvironment{Shaded}{\begin{snugshade}}{\end{snugshade}}
\newcommand{\AlertTok}[1]{\textcolor[rgb]{0.68,0.00,0.00}{#1}}
\newcommand{\AnnotationTok}[1]{\textcolor[rgb]{0.37,0.37,0.37}{#1}}
\newcommand{\AttributeTok}[1]{\textcolor[rgb]{0.40,0.45,0.13}{#1}}
\newcommand{\BaseNTok}[1]{\textcolor[rgb]{0.68,0.00,0.00}{#1}}
\newcommand{\BuiltInTok}[1]{\textcolor[rgb]{0.00,0.23,0.31}{#1}}
\newcommand{\CharTok}[1]{\textcolor[rgb]{0.13,0.47,0.30}{#1}}
\newcommand{\CommentTok}[1]{\textcolor[rgb]{0.37,0.37,0.37}{#1}}
\newcommand{\CommentVarTok}[1]{\textcolor[rgb]{0.37,0.37,0.37}{\textit{#1}}}
\newcommand{\ConstantTok}[1]{\textcolor[rgb]{0.56,0.35,0.01}{#1}}
\newcommand{\ControlFlowTok}[1]{\textcolor[rgb]{0.00,0.23,0.31}{\textbf{#1}}}
\newcommand{\DataTypeTok}[1]{\textcolor[rgb]{0.68,0.00,0.00}{#1}}
\newcommand{\DecValTok}[1]{\textcolor[rgb]{0.68,0.00,0.00}{#1}}
\newcommand{\DocumentationTok}[1]{\textcolor[rgb]{0.37,0.37,0.37}{\textit{#1}}}
\newcommand{\ErrorTok}[1]{\textcolor[rgb]{0.68,0.00,0.00}{#1}}
\newcommand{\ExtensionTok}[1]{\textcolor[rgb]{0.00,0.23,0.31}{#1}}
\newcommand{\FloatTok}[1]{\textcolor[rgb]{0.68,0.00,0.00}{#1}}
\newcommand{\FunctionTok}[1]{\textcolor[rgb]{0.28,0.35,0.67}{#1}}
\newcommand{\ImportTok}[1]{\textcolor[rgb]{0.00,0.46,0.62}{#1}}
\newcommand{\InformationTok}[1]{\textcolor[rgb]{0.37,0.37,0.37}{#1}}
\newcommand{\KeywordTok}[1]{\textcolor[rgb]{0.00,0.23,0.31}{\textbf{#1}}}
\newcommand{\NormalTok}[1]{\textcolor[rgb]{0.00,0.23,0.31}{#1}}
\newcommand{\OperatorTok}[1]{\textcolor[rgb]{0.37,0.37,0.37}{#1}}
\newcommand{\OtherTok}[1]{\textcolor[rgb]{0.00,0.23,0.31}{#1}}
\newcommand{\PreprocessorTok}[1]{\textcolor[rgb]{0.68,0.00,0.00}{#1}}
\newcommand{\RegionMarkerTok}[1]{\textcolor[rgb]{0.00,0.23,0.31}{#1}}
\newcommand{\SpecialCharTok}[1]{\textcolor[rgb]{0.37,0.37,0.37}{#1}}
\newcommand{\SpecialStringTok}[1]{\textcolor[rgb]{0.13,0.47,0.30}{#1}}
\newcommand{\StringTok}[1]{\textcolor[rgb]{0.13,0.47,0.30}{#1}}
\newcommand{\VariableTok}[1]{\textcolor[rgb]{0.07,0.07,0.07}{#1}}
\newcommand{\VerbatimStringTok}[1]{\textcolor[rgb]{0.13,0.47,0.30}{#1}}
\newcommand{\WarningTok}[1]{\textcolor[rgb]{0.37,0.37,0.37}{\textit{#1}}}

\usepackage{longtable,booktabs,array}
\usepackage{calc} % for calculating minipage widths
% Correct order of tables after \paragraph or \subparagraph
\usepackage{etoolbox}
\makeatletter
\patchcmd\longtable{\par}{\if@noskipsec\mbox{}\fi\par}{}{}
\makeatother
% Allow footnotes in longtable head/foot
\IfFileExists{footnotehyper.sty}{\usepackage{footnotehyper}}{\usepackage{footnote}}
\makesavenoteenv{longtable}
\usepackage{graphicx}
\makeatletter
\newsavebox\pandoc@box
\newcommand*\pandocbounded[1]{% scales image to fit in text height/width
  \sbox\pandoc@box{#1}%
  \Gscale@div\@tempa{\textheight}{\dimexpr\ht\pandoc@box+\dp\pandoc@box\relax}%
  \Gscale@div\@tempb{\linewidth}{\wd\pandoc@box}%
  \ifdim\@tempb\p@<\@tempa\p@\let\@tempa\@tempb\fi% select the smaller of both
  \ifdim\@tempa\p@<\p@\scalebox{\@tempa}{\usebox\pandoc@box}%
  \else\usebox{\pandoc@box}%
  \fi%
}
% Set default figure placement to htbp
\def\fps@figure{htbp}
\makeatother





\setlength{\emergencystretch}{3em} % prevent overfull lines

\providecommand{\tightlist}{%
  \setlength{\itemsep}{0pt}\setlength{\parskip}{0pt}}



 


\usepackage[document]{ragged2e}
\KOMAoption{captions}{tableheading}
\makeatletter
\@ifpackageloaded{caption}{}{\usepackage{caption}}
\AtBeginDocument{%
\ifdefined\contentsname
  \renewcommand*\contentsname{Table of contents}
\else
  \newcommand\contentsname{Table of contents}
\fi
\ifdefined\listfigurename
  \renewcommand*\listfigurename{List of Figures}
\else
  \newcommand\listfigurename{List of Figures}
\fi
\ifdefined\listtablename
  \renewcommand*\listtablename{List of Tables}
\else
  \newcommand\listtablename{List of Tables}
\fi
\ifdefined\figurename
  \renewcommand*\figurename{Figure}
\else
  \newcommand\figurename{Figure}
\fi
\ifdefined\tablename
  \renewcommand*\tablename{Table}
\else
  \newcommand\tablename{Table}
\fi
}
\@ifpackageloaded{float}{}{\usepackage{float}}
\floatstyle{ruled}
\@ifundefined{c@chapter}{\newfloat{codelisting}{h}{lop}}{\newfloat{codelisting}{h}{lop}[chapter]}
\floatname{codelisting}{Listing}
\newcommand*\listoflistings{\listof{codelisting}{List of Listings}}
\makeatother
\makeatletter
\makeatother
\makeatletter
\@ifpackageloaded{caption}{}{\usepackage{caption}}
\@ifpackageloaded{subcaption}{}{\usepackage{subcaption}}
\makeatother
\usepackage{bookmark}
\IfFileExists{xurl.sty}{\usepackage{xurl}}{} % add URL line breaks if available
\urlstyle{same}
\hypersetup{
  pdftitle={Arbeitsblatt -- Tupel und Listen in Python},
  colorlinks=true,
  linkcolor={blue},
  filecolor={Maroon},
  citecolor={Blue},
  urlcolor={Blue},
  pdfcreator={LaTeX via pandoc}}


\title{Arbeitsblatt -- Tupel und Listen in Python}
\author{}
\date{}
\begin{document}
\maketitle


\subsubsection{Aufgabe 1}\label{aufgabe-1}

Gegeben ist folgendes Tupel:

\begin{Shaded}
\begin{Highlighting}[numbers=left,,]
\NormalTok{primes }\OperatorTok{=}\NormalTok{ (}\DecValTok{2}\NormalTok{, }\DecValTok{3}\NormalTok{, }\DecValTok{5}\NormalTok{, }\DecValTok{7}\NormalTok{, }\DecValTok{11}\NormalTok{, }\DecValTok{13}\NormalTok{)}
\end{Highlighting}
\end{Shaded}

\begin{enumerate}
\def\labelenumi{\arabic{enumi}.}
\tightlist
\item
  Gib das \textbf{zweite Element} des Tupels aus.\\
\item
  Gib das \textbf{letzte Element} aus (nutze einen negativen Index).\\
\item
  Gib die \textbf{Länge des Tupels} aus.
\end{enumerate}

\begin{Shaded}
\begin{Highlighting}[numbers=left,,]
\CommentTok{\# Lösung zu Aufgabe 1}
\end{Highlighting}
\end{Shaded}

\subsubsection{Aufgabe 2}\label{aufgabe-2}

Erstelle eine Liste:

\begin{Shaded}
\begin{Highlighting}[numbers=left,,]
\NormalTok{numbers }\OperatorTok{=}\NormalTok{ [}\DecValTok{4}\NormalTok{, }\DecValTok{8}\NormalTok{, }\DecValTok{15}\NormalTok{, }\DecValTok{16}\NormalTok{, }\DecValTok{23}\NormalTok{, }\DecValTok{42}\NormalTok{]}
\end{Highlighting}
\end{Shaded}

\begin{enumerate}
\def\labelenumi{\arabic{enumi}.}
\tightlist
\item
  Gib das \textbf{erste} und das \textbf{letzte Element} aus.\\
\item
  Berechne die \textbf{Summe der Liste}.\\
\item
  Gib das \textbf{Minimum} und das \textbf{Maximum} der Liste aus.
\end{enumerate}

\begin{Shaded}
\begin{Highlighting}[numbers=left,,]
\CommentTok{\# Lösung zu Aufgabe 2}
\end{Highlighting}
\end{Shaded}

\subsubsection{Aufgabe 3}\label{aufgabe-3}

Gegeben:

\begin{Shaded}
\begin{Highlighting}[numbers=left,,]
\NormalTok{a }\OperatorTok{=}\NormalTok{ [}\DecValTok{1}\NormalTok{, }\DecValTok{2}\NormalTok{, }\DecValTok{3}\NormalTok{]}
\NormalTok{b }\OperatorTok{=}\NormalTok{ [}\DecValTok{4}\NormalTok{, }\DecValTok{5}\NormalTok{]}
\end{Highlighting}
\end{Shaded}

\begin{enumerate}
\def\labelenumi{\arabic{enumi}.}
\tightlist
\item
  Verbinde beide Listen zu einer neuen Liste.\\
\item
  Wiederhole die Liste \texttt{a} \textbf{dreimal}.\\
\item
  Verbinde das Tupel \texttt{(7,\ 8)} mit sich selbst.
\end{enumerate}

\begin{Shaded}
\begin{Highlighting}[numbers=left,,]
\CommentTok{\# Lösung zu Aufgabe 3}
\end{Highlighting}
\end{Shaded}

\subsubsection{Aufgabe 4}\label{aufgabe-4}

Gegeben:

\begin{Shaded}
\begin{Highlighting}[numbers=left,,]
\NormalTok{food }\OperatorTok{=}\NormalTok{ (}\StringTok{"Pizza"}\NormalTok{, }\StringTok{"Salat"}\NormalTok{, }\StringTok{"Burger"}\NormalTok{, }\StringTok{"Pasta"}\NormalTok{)}
\end{Highlighting}
\end{Shaded}

\begin{enumerate}
\def\labelenumi{\arabic{enumi}.}
\tightlist
\item
  Prüfe, ob \texttt{"Pizza"} im Tupel enthalten ist.\\
\item
  Prüfe, ob \texttt{"Sushi"} enthalten ist.\\
\item
  Gib eine passende Meldung aus, wenn \texttt{"Salat"} vorkommt.
\end{enumerate}

Nutze dafür den \textbf{\texttt{in}-Operator}.

\begin{Shaded}
\begin{Highlighting}[numbers=left,,]
\CommentTok{\# Lösung zu Aufgabe 4}
\end{Highlighting}
\end{Shaded}

\subsubsection{Aufgabe 5}\label{aufgabe-5}

Gegeben ist das Tupel:

\begin{Shaded}
\begin{Highlighting}[numbers=left,,]
\NormalTok{numbers }\OperatorTok{=}\NormalTok{ (}\DecValTok{10}\NormalTok{, }\DecValTok{20}\NormalTok{, }\DecValTok{30}\NormalTok{, }\DecValTok{40}\NormalTok{, }\DecValTok{50}\NormalTok{, }\DecValTok{60}\NormalTok{, }\DecValTok{70}\NormalTok{)}
\end{Highlighting}
\end{Shaded}

\begin{enumerate}
\def\labelenumi{\arabic{enumi}.}
\tightlist
\item
  Gib die \textbf{ersten drei Elemente} aus.\\
\item
  Gib die \textbf{letzten zwei Elemente} aus.\\
\item
  Gib eine \textbf{Kopie des gesamten Tupels} mithilfe von Slicing aus.
\end{enumerate}

\begin{Shaded}
\begin{Highlighting}[numbers=left,,]
\CommentTok{\# Lösung zu Aufgabe 5}
\end{Highlighting}
\end{Shaded}

\subsubsection{Aufgabe 6}\label{aufgabe-6}

Gegeben:

\begin{Shaded}
\begin{Highlighting}[numbers=left,,]
\NormalTok{point }\OperatorTok{=}\NormalTok{ (}\DecValTok{4}\NormalTok{, }\DecValTok{7}\NormalTok{)}
\end{Highlighting}
\end{Shaded}

\begin{enumerate}
\def\labelenumi{\arabic{enumi}.}
\tightlist
\item
  Weise die Werte des Tupels den Variablen \texttt{x} und \texttt{y}
  zu.\\
\item
  Gib beide Werte aus.
\end{enumerate}

Erweiterung:

\begin{Shaded}
\begin{Highlighting}[numbers=left,,]
\NormalTok{data }\OperatorTok{=}\NormalTok{ (}\DecValTok{42}\NormalTok{, (}\DecValTok{6}\NormalTok{, }\DecValTok{9}\NormalTok{), }\StringTok{"do"}\NormalTok{)}
\end{Highlighting}
\end{Shaded}

Entpacke das Tupel so, dass du auf die einzelnen Werte zugreifen kannst.

\begin{Shaded}
\begin{Highlighting}[numbers=left,,]
\CommentTok{\# Lösung zu Aufgabe 6}
\end{Highlighting}
\end{Shaded}

\subsubsection{Aufgabe 7}\label{aufgabe-7}

Gegeben:

\begin{Shaded}
\begin{Highlighting}[numbers=left,,]
\NormalTok{t }\OperatorTok{=}\NormalTok{ (}\DecValTok{3}\NormalTok{, }\DecValTok{6}\NormalTok{, }\DecValTok{9}\NormalTok{)}
\end{Highlighting}
\end{Shaded}

\begin{enumerate}
\def\labelenumi{\arabic{enumi}.}
\tightlist
\item
  Wandle das Tupel in eine \textbf{Liste} um.\\
\item
  Füge der Liste die Zahl \textbf{12} hinzu.\\
\item
  Wandle die Liste anschließend wieder in ein \textbf{Tupel} um.
\end{enumerate}

Nutze dafür \texttt{list()} und \texttt{tuple()}.

\begin{Shaded}
\begin{Highlighting}[numbers=left,,]
\CommentTok{\# Lösung zu Aufgabe 7}
\end{Highlighting}
\end{Shaded}

\subsubsection{Aufgabe 8}\label{aufgabe-8}

Gegeben ist die Liste:

\begin{Shaded}
\begin{Highlighting}[numbers=left,,]
\NormalTok{values }\OperatorTok{=}\NormalTok{ [}\DecValTok{12}\NormalTok{, }\DecValTok{5}\NormalTok{, }\DecValTok{8}\NormalTok{, }\DecValTok{21}\NormalTok{, }\DecValTok{3}\NormalTok{, }\DecValTok{17}\NormalTok{]}
\end{Highlighting}
\end{Shaded}

\begin{enumerate}
\def\labelenumi{\arabic{enumi}.}
\tightlist
\item
  Erzeuge eine \textbf{sortierte Version der Liste}.\\
\item
  Prüfe mit \texttt{any()}, ob eine Zahl \textbf{größer als 20} ist.\\
\item
  Prüfe mit \texttt{all()}, ob \textbf{alle Zahlen positiv} sind.
\end{enumerate}

\begin{Shaded}
\begin{Highlighting}[numbers=left,,]
\CommentTok{\# Lösung zu Aufgabe 8}
\end{Highlighting}
\end{Shaded}

\subsubsection{Zusatzaufgabe}\label{zusatzaufgabe}

Gegeben:

\begin{Shaded}
\begin{Highlighting}[numbers=left,,]
\NormalTok{letters }\OperatorTok{=}\NormalTok{ [}\StringTok{"p"}\NormalTok{, }\StringTok{"y"}\NormalTok{, }\StringTok{"t"}\NormalTok{, }\StringTok{"h"}\NormalTok{, }\StringTok{"o"}\NormalTok{, }\StringTok{"n"}\NormalTok{]}
\end{Highlighting}
\end{Shaded}

\begin{enumerate}
\def\labelenumi{\arabic{enumi}.}
\tightlist
\item
  Verbinde die Liste zu einem String mit \texttt{"".join()}.\\
\item
  Verbinde die Buchstaben mit \texttt{"-"}.
\end{enumerate}

Beispiel:

\begin{Shaded}
\begin{Highlighting}[numbers=left,,]
\NormalTok{p}\OperatorTok{{-}}\NormalTok{y}\OperatorTok{{-}}\NormalTok{t}\OperatorTok{{-}}\NormalTok{h}\OperatorTok{{-}}\NormalTok{o}\OperatorTok{{-}}\NormalTok{n}
\end{Highlighting}
\end{Shaded}

\begin{Shaded}
\begin{Highlighting}[numbers=left,,]
\CommentTok{\# Lösung zur Zusatzaufgabe}
\end{Highlighting}
\end{Shaded}





\end{document}
